\documentclass[10pt]{article}
\usepackage[margin=1.8cm]{geometry}
\usepackage[utf8]{inputenc}
\usepackage[T1]{fontenc}
\usepackage[spanish]{babel}
\shorthandoff{<>}
\usepackage{amsmath,booktabs,array}
\usepackage{xcolor}
\usepackage{tikz}
\usetikzlibrary{arrows.meta}
\setlength{\parindent}{0pt}
\newcommand{\campo}[1]{\underline{\hspace{#1}}}
\begin{document}

\begin{center}
{\LARGE Reporte de laboratorio: clasificación tabular con MLP}\\[0.4em]
Nombre: \campo{6cm}\hfill Fecha: \campo{3cm}\\
Repositorio o archivo entregado: \campo{10cm}
\end{center}

\textbf{Propósito.} Documentar una comparación reproducible entre un MLP convencional y un MLP residual con Batch Normalization, usando una división estratificada 800/200 y un conjunto de entrenamiento desbalanceado. Los valores mostrados como ejemplo son ficticios y deben sustituirse por resultados propios.

\section*{1. Datos, división y control de fuga}
\begin{tabular}{@{}llll@{}}
\toprule
Elemento & Valor & Evidencia o justificación & Estado\\
\midrule
Archivo público & train\_1000\_desbalanceado.csv & 1,000 ejemplos & Ejemplo\\
División & 800 train / 200 dev & Estratificada, semilla 31 & Ejemplo\\
Clases en train & 600 clase 0 / 200 clase 1 & Proporción conservada & Ejemplo\\
Estandarización & Media y desviación de train & Sin usar dev ni test & Ejemplo\\
Test privado & No usado durante selección & Sólo evaluación final & Ejemplo\\
\bottomrule
\end{tabular}

\vspace{0.7em}
\textbf{Respuesta.} Describa una fuente posible de fuga de información y cómo la evitó:

\vspace{2.0cm}

\section*{2. Estrategia ante el desbalance}
\begin{tabular}{@{}p{3.1cm}p{5.4cm}p{5.3cm}@{}}
\toprule
Decisión & Valor elegido & Justificación\\
\midrule
Pérdida & BCE ponderada; peso positivo = 3.0 & Compensa la relación 600/200 de train.\\
Selección de modelo & Costo esperado en dev & Falso negativo cuesta 5; falso positivo cuesta 1.\\
Métricas reportadas & Precision, recall, F1, especificidad y accuracy & Accuracy sola oculta el desempeño de la clase minoritaria.\\
Umbral & Se elige sólo en dev & Se congela antes del test privado.\\
\bottomrule
\end{tabular}

\vspace{0.6em}
\textbf{Respuesta.} Declare los costos y explique por qué el umbral 0.5 puede no ser adecuado:

\vspace{1.5cm}

\section*{3. Búsqueda de hiperparámetros}
\small
\begin{tabular}{@{}lcccccccc@{}}
\toprule
ID & Arquitectura & Capas/ancho & LR & Épocas & Costo dev & F1 dev & Recall dev & Umbral\\
\midrule
A & MLP & [24,16] & 0.003 & 100 & 43 & 0.81 & 0.90 & 0.41\\
B & MLP & [48,24] & 0.002 & 100 & 38 & 0.84 & 0.88 & 0.38\\
C & Residual + BN & 24 & 0.002 & 100 & 46 & 0.79 & 0.92 & 0.35\\
D & Residual + BN & 40 & 0.0015 & 100 & 41 & 0.82 & 0.90 & 0.40\\
\bottomrule
\end{tabular}
\normalsize

\vspace{0.5em}
\textbf{Decisión del ejemplo.} Se elige B por costo mínimo en desarrollo. Documente su regla de desempate y su modelo elegido:

\vspace{1.5cm}

\section*{4. Curvas de entrenamiento y desarrollo}
\begin{center}
\begin{tikzpicture}[x=0.85cm,y=3.2cm]
  \draw (0,0) -- (7.5,0) node[right]{Época};
  \draw (0,0) -- (0,1.25) node[above]{Pérdida};
  \foreach \x/\e in {1/20,2/40,3/60,4/80,5/100,6/120,7/140} \draw (\x,0) -- (\x,-.03) node[below,font=\scriptsize]{\e};
  \foreach \y/\e in {.25/.25,.5/.50,.75/.75,1/1.00} \draw (0,\y) -- (-.05,\y) node[left,font=\scriptsize]{\e};
  \draw[blue,thick] plot coordinates {(0,.95)(1,.61)(2,.43)(3,.31)(4,.25)(5,.21)(6,.18)(7,.16)};
  \draw[red,thick,dashed] plot coordinates {(0,.98)(1,.66)(2,.49)(3,.39)(4,.35)(5,.34)(6,.35)(7,.37)};
  \node[blue] at (5.9,.24) {\scriptsize train};
  \node[red] at (6.3,.43) {\scriptsize dev};
\end{tikzpicture}
\hspace{1cm}
\begin{tikzpicture}[x=0.85cm,y=3.2cm]
  \draw (0,0) -- (7.5,0) node[right]{Época};
  \draw (0,0) -- (0,1.25) node[above]{Métrica};
  \foreach \x/\e in {1/20,2/40,3/60,4/80,5/100,6/120,7/140} \draw (\x,0) -- (\x,-.03) node[below,font=\scriptsize]{\e};
  \foreach \y/\e in {.25/.25,.5/.50,.75/.75,1/1.00} \draw (0,\y) -- (-.05,\y) node[left,font=\scriptsize]{\e};
  \draw[blue,thick] plot coordinates {(0,.58)(1,.70)(2,.77)(3,.82)(4,.85)(5,.86)(6,.85)(7,.84)};
  \draw[green!60!black,thick,dashed] plot coordinates {(0,.42)(1,.60)(2,.71)(3,.78)(4,.82)(5,.84)(6,.83)(7,.82)};
  \node[blue] at (5.8,.93) {\scriptsize accuracy dev};
  \node[green!60!black] at (5.4,.74) {\scriptsize F1 dev};
\end{tikzpicture}
\end{center}

\textbf{Interpretación.} Indique la época o checkpoint seleccionado y explique si observa sobreajuste, subajuste o estabilidad:

\vspace{1.4cm}

\section*{5. Selección de umbral y reporte de métricas}
\begin{center}
\begin{tikzpicture}[x=8cm,y=.07cm]
  \draw (0,0) -- (1.08,0) node[right]{Umbral};
  \draw (0,0) -- (0,52) node[above]{Costo dev};
  \draw[gray] (.1,0)--(.1,50) (.3,0)--(.3,50) (.5,0)--(.5,50) (.7,0)--(.7,50) (.9,0)--(.9,50);
  \draw[red,thick] plot coordinates {(.05,49)(.15,38)(.25,29)(.35,21)(.40,18)(.45,20)(.55,27)(.65,39)(.75,48)(.85,51)};
  \draw[dashed] (.40,0)--(.40,18);
  \node[below] at (.40,0) {\scriptsize 0.40};
\end{tikzpicture}
\end{center}

\begin{center}
\begin{tabular}{@{}lccccc@{}}
\toprule
Conjunto & Accuracy & Precision & Recall & F1 & Especificidad\\
\midrule
Desarrollo, umbral 0.40 & 0.91 & 0.78 & 0.92 & 0.84 & 0.90\\
Test privado, umbral congelado & \campo{1cm} & \campo{1cm} & \campo{1cm} & \campo{1cm} & \campo{1cm}\\
\bottomrule
\end{tabular}
\end{center}

\textbf{Matriz de confusión en desarrollo (ejemplo ficticio).}
\[
\begin{array}{c|cc}
 & \text{Pred. 0} & \text{Pred. 1}\\ \hline
\text{Real 0} & 135 & 15\\
\text{Real 1} & 4 & 46
\end{array}
\]

\section*{6. Conclusión reproducible}
\begin{enumerate}
\item Arquitectura final, pesos y ruta del artefacto: \campo{10cm}
\item Umbral congelado y costos asumidos: \campo{10cm}
\item Cambio que haría antes de desplegar: \campo{10cm}
\end{enumerate}
\end{document}
